% GENERATED FILE. Edit canonical lessons, then run npm run generate:books.
% canonical-source-hash: fnv1a64:28a89918c2e8d603
% canonical-lessons: PA-C126-ghah, PA-C126-bij, PA-C126-jarh, PA-C126-kali, PA-C126-tapu

\chapter{Grass, Seed, Root, Bud, Island}
\label{ch:pa-grass-seed-root-bud-island}
\hlchaptermodality{126}

\begin{chapteropening}
\textbf{By the end of this chapter:} \emph{I can say grass, a seed, a root, a bud and an island.}

Complete the last lesson of chapter 126: I can say grass, a seed, a root, a bud and an island.
\end{chapteropening}

\section[\texorpdfstring{\pa{ਘਾਹ} (ghāh)}{ਘਾਹ (ghāh)}]{\pa{ਘਾਹ} (ghāh) — grass}

\label{lesson:PA-C126-ghah}



\begin{quote}
Before the new one: say the Punjabi for cumin, then the Punjabi for a fritter.
\end{quote}

\subsection*{You'll want to know: \pa{ਘਾਹ}}
\textbf{\pa{ਘਾਹ}} — \emph{ghāh} — \textquotedblleft{}grass\textquotedblright{}.

\begin{cousinweb}
\textbf{\pa{ਘਾਹ}} (\emph{ghāh}) is a tadbhava, a word inherited from Sanskrit and reshaped by sound change. It comes from \emph{ghāsa}, \textquotedblleft{}food, pasture grass\textquotedblright{} (CDIAL 4471).
\end{cousinweb}

\begin{grammarlens}[title={What you've built}]
One more word for this chapter.
\end{grammarlens}

\subsection*{Your turn}
\begin{itemize}
\raggedright
  \item \emph{Say it:} \emph{ghāh}
  \item \emph{Say it:} \emph{ghāh}, once more
  \item \emph{Say it:} say \emph{ghāh}
  \item \emph{Recall:} say the Punjabi for cumin, then the Punjabi for a fritter, then say \emph{ghāh} again
\end{itemize}

\begin{tcolorbox}[breakable,colback=teal!4,colframe=teal!35!black,title={Before you move on}]
What does \pa{ਘਾਹ} mean? (\textquotedblleft{}Grass\textquotedblright{}.) Say it once more.
\end{tcolorbox}
\section[\texorpdfstring{\pa{ਬੀਜ} (bīj)}{ਬੀਜ (bīj)}]{\pa{ਬੀਜ} (bīj) — a seed}

\label{lesson:PA-C126-bij}



\begin{quote}
Before the new one: say the Punjabi for a fritter, then the Punjabi for grass.
\end{quote}

\subsection*{You'll want to know: \pa{ਬੀਜ}}
\textbf{\pa{ਬੀਜ}} — \emph{bīj} — \textquotedblleft{}a seed\textquotedblright{}.

\begin{cousinweb}
\textbf{\pa{ਬੀਜ}} (\emph{bīj}) is a semi-tatsama: Sanskrit \emph{bīja}, \textquotedblleft{}seed, semen\textquotedblright{}, borrowed after the old sound changes had run their course and then lightly reshaped (CDIAL 9250). CDIAL lists the inherited Punjabi form as \emph{bī}.
\end{cousinweb}

\begin{grammarlens}[title={What you've built}]
One more word for this chapter.
\end{grammarlens}

\subsection*{Your turn}
\begin{itemize}
\raggedright
  \item \emph{Say it:} \emph{bīj}
  \item \emph{Say it:} \emph{bīj}, once more
  \item \emph{Say it:} say \emph{bīj}
  \item \emph{Recall:} say the Punjabi for a fritter, then the Punjabi for grass, then say \emph{bīj} again
\end{itemize}

\begin{tcolorbox}[breakable,colback=teal!4,colframe=teal!35!black,title={Before you move on}]
What does \pa{ਬੀਜ} mean? (\textquotedblleft{}A seed\textquotedblright{}.) Say it once more.
\end{tcolorbox}
\section[\texorpdfstring{\pa{ਜੜ੍ਹ} (jaṛh)}{ਜੜ੍ਹ (jaṛh)}]{\pa{ਜੜ੍ਹ} (jaṛh) — a root}

\label{lesson:PA-C126-jarh}



\begin{quote}
Before the new one: say the Punjabi for grass, then the Punjabi for a seed.
\end{quote}

\subsection*{You'll want to know: \pa{ਜੜ੍ਹ}}
\textbf{\pa{ਜੜ੍ਹ}} — \emph{jaṛh} — \textquotedblleft{}a root\textquotedblright{}.

\begin{cousinweb}
\textbf{\pa{ਜੜ੍ਹ}} (\emph{jaṛh}) is a tadbhava, a word inherited from Sanskrit and reshaped by sound change. It comes from \emph{jaṭā}, \textquotedblleft{}hair twisted together; fibrous root\textquotedblright{} (CDIAL 5086). CDIAL reports views that the Sanskrit word is from Dravidian, or ultimately from Munda.
\end{cousinweb}

\begin{grammarlens}[title={What you've built}]
One more word for this chapter.
\end{grammarlens}

\subsection*{Your turn}
\begin{itemize}
\raggedright
  \item \emph{Say it:} \emph{jaṛh}
  \item \emph{Say it:} \emph{jaṛh}, once more
  \item \emph{Say it:} say \emph{jaṛh}
  \item \emph{Recall:} say the Punjabi for grass, then the Punjabi for a seed, then say \emph{jaṛh} again
\end{itemize}

\begin{tcolorbox}[breakable,colback=teal!4,colframe=teal!35!black,title={Before you move on}]
What does \pa{ਜੜ੍ਹ} mean? (\textquotedblleft{}A root\textquotedblright{}.) Say it once more.
\end{tcolorbox}
\section[\texorpdfstring{\pa{ਕਲੀ} (kalī)}{ਕਲੀ (kalī)}]{\pa{ਕਲੀ} (kalī) — a bud}

\label{lesson:PA-C126-kali}



\begin{quote}
Before the new one: say the Punjabi for a seed, then the Punjabi for a root.
\end{quote}

\subsection*{You'll want to know: \pa{ਕਲੀ}}
\textbf{\pa{ਕਲੀ}} — \emph{kalī} — \textquotedblleft{}a bud\textquotedblright{}.

\begin{cousinweb}
\textbf{\pa{ਕਲੀ}} (\emph{kalī}) is a tadbhava, a word inherited from Sanskrit and reshaped by sound change. It comes from \emph{kali}, \textquotedblleft{}unblown bud\textquotedblright{} (CDIAL 2934).
\end{cousinweb}

\begin{grammarlens}[title={What you've built}]
One more word for this chapter.
\end{grammarlens}

\subsection*{Your turn}
\begin{itemize}
\raggedright
  \item \emph{Say it:} \emph{kalī}
  \item \emph{Say it:} \emph{kalī}, once more
  \item \emph{Say it:} say \emph{kalī}
  \item \emph{Recall:} say the Punjabi for a seed, then the Punjabi for a root, then say \emph{kalī} again
\end{itemize}

\begin{tcolorbox}[breakable,colback=teal!4,colframe=teal!35!black,title={Before you move on}]
What does \pa{ਕਲੀ} mean? (\textquotedblleft{}A bud\textquotedblright{}.) Say it once more.
\end{tcolorbox}
\section[\texorpdfstring{\pa{ਟਾਪੂ} (ṭāpū)}{ਟਾਪੂ (ṭāpū)}]{\pa{ਟਾਪੂ} (ṭāpū) — an island}

\label{lesson:PA-C126-tapu}



\begin{quote}
Before the new one: say the Punjabi for a root, then the Punjabi for a bud.
\end{quote}

\subsection*{You'll want to know: \pa{ਟਾਪੂ}}
\textbf{\pa{ਟਾਪੂ}} — \emph{ṭāpū} — \textquotedblleft{}an island\textquotedblright{}.

\begin{cousinweb}
\textbf{\pa{ਟਾਪੂ}} (\emph{ṭāpū}) is a tadbhava, a word inherited and reshaped by sound change. It comes from \emph{*ṭāppuka}, \textquotedblleft{}island\textquotedblright{}, a form CDIAL reconstructs from its descendants (CDIAL 5456).
\end{cousinweb}

\begin{grammarlens}[title={What you've built}]
One more word for this chapter.
\end{grammarlens}

\subsection*{Your turn}
\begin{itemize}
\raggedright
  \item \emph{Say it:} \emph{ṭāpū}
  \item \emph{Say it:} \emph{ṭāpū}, once more
  \item \emph{Say it:} say \emph{ṭāpū}
  \item \emph{Recall:} say the Punjabi for a root, then the Punjabi for a bud, then say \emph{ṭāpū} again
\end{itemize}

\begin{tcolorbox}[breakable,colback=teal!4,colframe=teal!35!black,title={Before you move on}]
What does \pa{ਟਾਪੂ} mean? (\textquotedblleft{}An island\textquotedblright{}.) Say it once more.
\end{tcolorbox}
